\begin{figure*}[t]
    \centering
    \includegraphics[width = 0.9\linewidth]{figs/comparison_menglei.pdf}\\
    \setlength{\tabcolsep}{0\linewidth}
    \begin{tabular}{C{0.2191\linewidth}C{0.3154\linewidth}C{0.3154\linewidth}}
        {\footnotesize(a) learnable pooling} & {\footnotesize(b) pooling by rasterization} & {\footnotesize(c) pooling by spatial proximity} \\
    \end{tabular}\vspace{-5pt}
    \caption{\textbf{Issues of existing multi-level GNNs.} (a) A learnable pooling \citep{gao2019graph} may lead to loss of connectivity even after $2^{\text{nd}}$-order enhancement. (b) A pooling by rasterization \citep{lino2021simulating,lino2022towards,lino2022multi} and (c) by spatial proximity \citep{liu2021multi,fortunato2022multiscale} can lead to wrong connections across the boundaries at the coarser level.}\vspace{-10pt}
    \label{fig:compare}
\end{figure*}
\section{Introduction}
\begin{figure*}[t]
    \centering
    \includegraphics[width=0.9\linewidth]{figs/example_bistride_anot.pdf}\vspace{2pt}\\
    \setlength{\tabcolsep}{0\linewidth}
    \begin{tabular}{C{0.05\linewidth}C{0.125\linewidth}C{0.275\linewidth}C{0.2\linewidth}C{0.15\linewidth}}
         & {\scriptsize(a) \cylinder~} & {\scriptsize(b) \airfoil~} & {\scriptsize(c) \plate~} & {\scriptsize(d) \fonts~} \\
    \end{tabular}\vspace{-5pt}
    \caption{\textbf{Example multi-level graphs produced by bi-stride pooling.} Our datasets contain both Eulerian and Lagrangian systems. Many meshes are highly irregular and contain massive self-contact, challenging to build coarser-level connections by spatial proximity. Our bi-stride strategy only relies on topological information and is proven to be robust and reliable on arbitrary geometry.}
    \vspace{-10pt}
    \label{fig:example:bistride}
\end{figure*}
Simulating physical systems through numerically solving partial differential equations (PDEs) is a key area in science and engineering, with applications ranging from solid mechanics~\cite{jiang2016material,li2020incremental}, to fluid-~\cite{bridson2015fluid,cao2018liquid}, aerodynamics~\cite{cao2022efficient}, and heat transfer~\cite{cao2019heat}.
However, traditional numerical solvers are often computationally expensive, particularly for time-sensitive applications such as iterative design optimization, where fast online inference is desired.
In recent years, the machine learning community has shown great interest in improving efficiency or replacing traditional solvers with learned models. These works include both end-to-end frameworks~\cite{grzeszczuk1998neuroanimator,obiols2020cfdnet} and those utilizing physics-informed losses~\cite{raissi2019physics,karniadakis2021physics,sun2020surrogate}.
Many existing works apply convolutional neural networks (CNNs)~\cite{fukushima1982neocognitron} to physical systems residing on two- or three-dimensional structured grids~\cite{guo2016convolutional,tompson2017accelerating,kim2019deep,fotiadis2020comparing}. However, the strict dependency on regular domain shapes makes it non-trivial to be applied on unstructured meshes. While it is possible to deform simple irregular domains into rectangular shapes to apply CNNs~\cite{gao2021phygeonet,li2022fourier}, the challenge remains for domains with complex topologies, which are common in practice.

As a result, the use of graph neural networks (GNNs) in physics-based simulations on unstructured meshes has gained significant attention in recent years~\cite{battaglia2018relational,sanchez2018graph,belbute2020combining,pfaff2020learning,sanchez2020learning,harsch2021direct,gao2022physics}.
The naive GNN approach stacks multiple MPs to model information propagation throughout space. However, as the graph size increases, this approach faces two major challenges:
(1) \textbf{Complexity}: as both the number of nodes to be processed and the MP iterations increase linearly, a quadratic complexity becomes inevitable for both the running time and memory usage of the computational graph~\cite{fortunato2022multiscale}.
(2) \textbf{Oversmoothing}: the graph convolution can be seen as a low-pass filter that suppresses the higher-frequency signals~\cite{chen2020measuring, li2020multipole}. Then the stacked MPs iteratively project the information onto the eigenspace of the graph while all higher-frequency signals are smoothed out, making the training more difficult.

To address these limitations, researchers have begun introducing multi-scale GNNs (MS-GNNs) for physics-based simulation~\cite{li2020multipole,liu2021multi,lino2021simulating,fortunato2022multiscale,lino2022multi,lino2022towards}. The multi-scale approach mitigates over-smoothing by building sub-level graphs at coarser resolutions, resulting in longer-range interactions and fewer MP iterations.
The existing approaches for constructing the multi-scale structure include: utilizing spatial proximity to generate sub-level graphs at coarser levels~\cite{lino2021simulating,liu2021multi,lino2022towards}; applying Guillard's coarsening algorithm~\cite{guillard1993node}~\cite{lino2022multi}; manually drawing coarser meshes for the original geometry~\cite{liu2021multi,fortunato2022multiscale}; or randomly pooling nodes and applying factorization to the adjacency matrix~\cite{li2020multipole}. However, these solutions all have their limitations.
For example, learnable or random pooling can introduce artificial partitions in the sub-level graphs (Fig.~\ref{fig:compare}. (a)), even with adjacency enhancement, which impedes information exchange across partitions, spatial proximity can lead to wrong edges across the boundaries at coarser levels (Fig.~\ref{fig:compare}. (b) and (c)); Guillard's algorithm only applies to 2D triangle meshes; and manually drawing tens of thousands of meshes is too labor-intensive.
We observe that all these limitations originate from the unmatured operations: \textit{pooling} and \textit{building graph connections at coarser levels}. We desire to design operations that: 1) conserve the correct connections at coarser levels, 2) do not introduce edges that blur the boundaries, 3) are general for any mesh type, and 4) are automatic.

We tackle these challenges with two progressive contributions:
\begin{itemize}
    \item \textbf{First}, we introduce a novel yet simple pooling strategy, \textit{bi-stride}. Bi-stride is inspired by the bi-partition determination in DAG (directed acyclic graph). It pools all nodes on every other BFS (breadth-first-search) frontier, such that a $2^{\text{nd}}$-powered adjacency enhancement \revision{($\mA\leftarrow\mA^2$, where $\mA$ is the adjacency matrix of the graph)} conserves all the correct connectivity. Bi-stride solely uses the input mesh without the need for spatial proximity, is general for any mesh type, and is fully automatic.
    \item \textbf{Second}, bi-stride pooling conserves direct connections between pooled/un-pooled nodes; Utilizing this advantage, one MP suffices to exchange the information between pooled and unpooled nodes before moving into the adjacent level; We also design a non-parameterized aggregating and returning method, resembling the interpolation in \revision{U-Net}, to handle the transition between adjacent levels. These simplifications significantly reduce the computational requirements compared to SOTAs.
\end{itemize}
Together, these two contributions give birth to our \textit{Bi-Stride Multi-Scale GNN} (\textit{BSMS-GNN}), a novel framework representing a significant advancement in the field of learned mesh-based simulations, particularly for the deployment in real industrial applications where meshes are often complex in geometry and large in size.
